% Card: Lucena's position (rook endings)
\paragraph{Lucena's position}\index{Rook endings!Lucena's position}

\diagram[lucena]{8/6P1/3k4/6K1/6R1/8/8/6r1 b - - 25 82}

The pawn is one step from queening, but the white king first had to leave the queening square,
and the black rook checks it from behind. The idea is simple: White brings his rook to the fourth
rank. The king walks out towards its rook, and when it reaches the fifth rank the rook steps in
on g4 and blocks the checks. The black king is far away and cannot stop the pawn. This is
\vocab[rook endings]{lucena}{Lucena's position}{The attacker's pawn is on the seventh rank with its
king in front of it, and the defending king is cut off. The attacker wins by bringing his rook to
the fourth rank, walking his king out, and blocking the checks from behind with the rook}.
